\itemhead{复习题一 · 第5题（教材p.29）}{XJ-B1-C01-RV-P029-Y0994}
\label{XJ-B1-C01-RV-P029-Y0994}
{\footnotesize 湘教版教材p.29 / PDF第36页 / 复习题一 / 第5题}
\field{原题}{设集合 $A=\{2,3,a^2+4a+2\}$，集合 $B=\{0,7,a^2+4a-2,2-a\}$，这里 $a$ 是某个正数，且 $7\in A$，求集合 $B$。}
\field{参考答案}{$a^2+4a+2=7$ 得 $a=1$ 或 $-5$，由正数条件取1；$B=\{0,7,3,1\}$。}
{\small \field{关联知识}{S01 集合、元素与常用数集（教材帮p.1、2 / PDF10、11）}
\field{方法与易错点}{先辨认研究对象和数集范围，再代入定义判断；含参集合逐个来源分类。 易错：把自然数集误作正整数集；混淆元素和单元素集合；候选参数造成元素重复。}
\field{答案依据}{答案册相关题组：PDF第8、9页（册内页7、8）；按当前题库原题逐项重算，不按旧版题号机械对套。}}
